\chapter{Intervalle und konvexe Teilmengen}
\label{b15guide:ch:intervals}

\section{Intervallnotation und Zerlegungen}

\begin{definition}[Intervalle und Strahlen]
\label{b15guide:def:intervals}
Für $a,b\in A$ verwenden wir
\[
\begin{aligned}
 [a,b]&=\{x\in A\mid a\leq x\leq b\},&
 (a,b)&=\{x\in A\mid a<x<b\},\\
 [a,b)&=\{x\in A\mid a\leq x<b\},&
 (a,b]&=\{x\in A\mid a<x\leq b\}.
\end{aligned}
\]
Ferner schreiben wir
\[
\begin{aligned}
 A_{<a}&=\{x\in A\mid x<a\},&
 A_{\leq a}&=\{x\in A\mid x\leq a\},\\
 A_{>a}&=\{x\in A\mid a<x\},&
 A_{\geq a}&=\{x\in A\mid a\leq x\}.
\end{aligned}
\]
Das \emph{unorientierte abgeschlossene Intervall} zwischen
$a$ und $b$ ist
\[
 [a,b]_{\mathrm u}=[\min(a,b),\max(a,b)].
\]
\end{definition}

\begin{theorem}[Zerlegung am Punkt und am Intervall]
\label{b15guide:thm:partitions}
Für jedes $a\in A$ ist
\[
 A=A_{<a}\mathbin{\dot\cup}\{a\}\mathbin{\dot\cup}A_{>a}.
\]
Für $a\leq b$ gilt außerdem die disjunkte Zerlegung
\[
 A=A_{<a}\mathbin{\dot\cup}[a,b]\mathbin{\dot\cup}A_{>b}.
\]
Insbesondere sind $A\setminus A_{\leq a}=A_{>a}$
und $A\setminus A_{<a}=A_{\geq a}$.
\end{theorem}

\begin{theorem}[Schnitt zweier abgeschlossener Intervalle]
\label{b15guide:thm:interval-intersection}
Für $a,b,c,d\in A$ gilt
\[
 [a,b]\cap[c,d]=[\max(a,c),\min(b,d)].
\]
Dieser Schnitt ist genau dann nichtleer, wenn
$\max(a,c)\leq\min(b,d)$ gilt.
Sind $a\leq b$ und $c\leq d$, so vereinfacht sich dies zu
\[
 [a,b]\cap[c,d]\ne\varnothing
 \quad\Longleftrightarrow\quad(a\leq d\land c\leq b).
\]
\end{theorem}

\begin{theorem}[Zerlegung durch einen Zwischenpunkt]
\label{b15guide:thm:interval-split}
Für $a\leq b\leq c$ gilt
\[
 [a,c]=[a,b]\cup[b,c],\qquad [a,b]\cap[b,c]=\{b\}.
\]
Für $a<b<c$ gilt die disjunkte Zerlegung
\[
 (a,c)=(a,b)\mathbin{\dot\cup}\{b\}\mathbin{\dot\cup}(b,c).
\]
\end{theorem}

\section{Konvexität und Trennung}

\begin{definition}[Ordnungskonvexität und konvexe Hülle]
\label{b15guide:def:convex}
Eine Teilmenge $C\subseteq A$ heißt \emph{ordnungskonvex},
wenn für $a,b\in C$ und $x\in A$ aus $a\leq x\leq b$
stets $x\in C$ folgt.
Mit $\operatorname{conv}_A(S)$ bezeichnen wir die kleinste
ordnungskonvexe Teilmenge von $A$, die $S$ enthält.
Sie ist der Durchschnitt aller ordnungskonvexen Obermengen von $S$.
\end{definition}

\begin{theorem}[Konvexe Hülle einer endlichen Teilmenge]
\label{b15guide:thm:finite-hull}
Für jede nichtleere endliche Teilmenge $S\subseteq A$ gilt
\[
 \operatorname{conv}_A(S)=[\min S,\max S].
\]
Insbesondere ist $\operatorname{conv}_A(\{a,b\})=[a,b]_{\mathrm u}$.
\end{theorem}

\begin{theorem}[Trennung disjunkter konvexer Mengen]
\label{b15guide:thm:separation}
Sind $C,D\subseteq A$ nichtleer, ordnungskonvex und disjunkt,
so gilt genau eine der Aussagen
\[
 \forall c\in C\,\forall d\in D\;(c<d),\qquad
 \forall c\in C\,\forall d\in D\;(d<c).
\]
Die gesamten Mengen liegen somit in einer eindeutigen Reihenfolge.
\end{theorem}

\begin{theorem}[Vereinigung überlappender konvexer Mengen]
\label{b15guide:thm:convex-union}
Sind $C,D\subseteq A$ ordnungskonvex und ist
$C\cap D\ne\varnothing$, so ist auch $C\cup D$ ordnungskonvex.
\end{theorem}

\begin{theorem}[Endliche Helly-Eigenschaft für Ordnungsintervalle]
\label{b15guide:thm:helly}
Es sei $\mathcal C$ eine nichtleere endliche Familie
nichtleerer ordnungskonvexer Teilmengen von $A$.
Dann gilt
\[
 \bigcap_{C\in\mathcal C} C\ne\varnothing
 \quad\Longleftrightarrow\quad
 \forall C,D\in\mathcal C\;(C\cap D\ne\varnothing).
\]
Die Mengen brauchen keine Endpunkte in $A$ zu besitzen.
Eine Aussage über unendliche Familien ist damit nicht verbunden.
\end{theorem}

\section{Der mittlere Wert dreier Elemente}

\begin{definition}[Median]
\label{b15guide:def:median}
Für $x,y,z\in A$ heißt
\[
 \operatorname{med}(x,y,z)=
 \max\bigl(\min(x,y),\min(y,z),\min(z,x)\bigr)
\]
der \emph{Median} der drei Elemente. Bei drei Argumenten
bedeutet $\max$ das Maximum ihrer Menge; analog verwenden wir $\min$.
\end{definition}

\begin{theorem}[Auswahl und Charakterisierung des Medians]
\label{b15guide:thm:median}
Der Median liegt in $\{x,y,z\}$ und ist der mittlere Wert
einer nichtabsteigenden Anordnung der drei Argumente,
wobei Wiederholungen erhalten bleiben. Es gilt auch
\[
 \operatorname{med}(x,y,z)=
 \min\bigl(\max(x,y),\max(y,z),\max(z,x)\bigr).
\]
Er ist das einzige $m\in A$, das gleichzeitig zwischen
$x,y$, zwischen $y,z$ und zwischen $z,x$ liegt:
\[
\begin{aligned}
 \min(x,y)&\leq m\leq\max(x,y),\\
 \min(y,z)&\leq m\leq\max(y,z),\\
 \min(z,x)&\leq m\leq\max(z,x).
\end{aligned}
\]
\end{theorem}
